\documentclass[11pt]{article}
\usepackage[margin=1in]{geometry}
\usepackage{amsmath,amssymb,amsthm,booktabs,graphicx,mathrsfs}
\usepackage[hidelinks]{hyperref}
\newtheorem{theorem}{Theorem}
\newtheorem{lemma}{Lemma}
\newtheorem{corollary}{Corollary}
\newcommand{\C}{\mathbb C}
\newcommand{\Gfam}{\mathscr G}
\newcommand{\ip}[2]{\langle #1,#2\rangle}
\setlength{\parindent}{0pt}
\setlength{\parskip}{0.65em}

\title{Multiplicity Obstructs a Universal Orthogonal Decomposition:\\
An $S_3$ Counterexample to Conjecture 1.1}
\author{Run \texttt{fa\_banach\_001}, lane 2}
\date{11 August 2026}

\begin{document}
\maketitle

\textbf{Status.} Candidate counterexample, likely valid, for human review.

\section{Source conjecture}

Hokamp's Conjecture~1.1 asks whether, whenever a compact group $G$ acts
continuously and transitively on a compact Hausdorff space $X$, there is one
collection $\Gfam$ of closed, pairwise $L^2$-orthogonal, $G$-minimal
subspaces of $C(X)$ such that every uniformly closed $G$-invariant subspace
of $C(X)$ is the uniform closure of the direct sum of a subcollection of
$\Gfam$ \cite[Conjecture~1.1, PDF p.~1]{Hokamp}.

\begin{center}
\includegraphics[width=\textwidth]{figures/open_problem_crop.png}
\end{center}

Here $C(X)$ consists of complex-valued functions, and orthogonality is with
respect to the invariant probability measure on $X$.

\section{Claimed contribution}

\begin{theorem}[Exact obstruction]\label{thm:characterization}
Let a compact group $G$ act continuously and transitively on a compact
Hausdorff space $X$, and let $U$ be the resulting unitary representation on
$L^2(X)$.  The conclusion of Conjecture~1.1 holds for $(G,X)$ if and only if
$U$ is multiplicity-free.

Equivalently, after writing $X\simeq G/K$ for a stabilizer $K$, the conclusion
holds if and only if
\[
  \dim V_\pi^K\leq 1
\]
for every irreducible unitary representation $V_\pi$ of $G$.
\end{theorem}

\begin{corollary}[Finite counterexample]\label{cor:s3}
Conjecture~1.1 is false.  Take $G=S_3$, let $X=S_3$ with the discrete
topology, and let $G$ act on $X$ by left translation.  This is a continuous
transitive action of a compact group on a compact Hausdorff space, but no
collection $\Gfam$ having the asserted universal property exists.
\end{corollary}

\section{Proof intuition}

The conjecture asks a single orthogonal list of irreducible subspaces to
generate \emph{every} invariant subspace merely by selecting entries from
that list.  This is possible when every irreducible type occurs once: the
isotypic pieces are then themselves irreducible, and Peter--Weyl spectral
projections recover every invariant subspace from a subset of them.

Multiplicity makes the requested list impossible.  In two equivalent copies
$V\oplus V$, both the first coordinate $V\oplus0$ and the diagonal
$\{(v,v):v\in V\}$ are minimal invariant subspaces, but they are not
orthogonal.  The universal property would force \emph{every} minimal
invariant subspace to be an entry of the same pairwise-orthogonal list.  The
left regular representation of $S_3$ contains two copies of its standard
two-dimensional representation, so it supplies a six-point counterexample.

\section{Proof of the characterization}

We use the following immediate consequence of the word ``every'' in the
source conjecture.

\begin{lemma}[Minimal spaces must be listed]\label{lem:listed}
Suppose a collection $\Gfam$ has the universal property in Conjecture~1.1.
Then every nonzero closed $G$-minimal subspace $M\subset C(X)$ is itself a
member of $\Gfam$.
\end{lemma}

\begin{proof}
By the universal property, $M$ is the closure of the direct sum of a
subcollection $\mathscr S\subset\Gfam$.  The subcollection is nonempty, and
each $H\in\mathscr S$ is contained in that direct sum and hence in $M$.
Since $H$ is nonzero and $G$-invariant while $M$ is minimal, $H=M$.
Thus $M\in\Gfam$.
\end{proof}

\begin{proof}[Proof of Theorem~\ref{thm:characterization}]
Let $\mu$ be the invariant probability measure.  Peter--Weyl theory gives
the isotypic decomposition
\[
 L^2(X,\mu)=\widehat\bigoplus_{\pi\in\widehat G}
             V_\pi\otimes M_\pi,
\]
where $M_\pi$ is the multiplicity space \cite{PeterWeyl,Folland}.  Because
$X\simeq G/K$, Frobenius reciprocity identifies
$M_\pi$ with $(V_\pi^*)^K$, so
$\dim M_\pi=\dim V_\pi^K<\infty$.

The finite-dimensional isotypic space
$H_\pi=V_\pi\otimes M_\pi$ consists of continuous functions.  One direct
way to see this is to use the character projection
\[
 P_\pi f=d_\pi\int_G \overline{\chi_\pi(g)}\,U_gf\,dg.
\]
It maps $C(X)$ to $C(X)$ and is the $L^2$-orthogonal projection onto
$H_\pi$.  Since $C(X)$ is dense in $L^2(X)$ and $H_\pi$ is finite
dimensional, $P_\pi(C(X))=H_\pi$.

Assume first that some $\dim M_\pi\geq2$.  Choose orthonormal vectors
$e_1,e_2\in M_\pi$.  Inside $C(X)$ consider
\[
 H_1=V_\pi\otimes\C e_1,
 \qquad
 H_\Delta=V_\pi\otimes\C(e_1+e_2).
\]
Both are closed $G$-minimal subspaces.  They are distinct and not
orthogonal, since for $0\neq v\in V_\pi$,
\[
 \ip{v\otimes e_1}{v\otimes(e_1+e_2)}=\lVert v\rVert^2\neq0.
\]
If the conjectured family existed, Lemma~\ref{lem:listed} would put both
$H_1$ and $H_\Delta$ into one pairwise-orthogonal family, a contradiction.
Thus the conjecture implies multiplicity-freeness.

Conversely, suppose $\dim M_\pi\leq1$ for every $\pi$.  Each nonzero
$H_\pi$ is then $G$-minimal, and the collection of these $H_\pi$ is pairwise
orthogonal.  Let $Y\subset C(X)$ be uniformly closed and $G$-invariant.
The Bochner-integral formula above shows that $P_\pi Y\subset Y$.
Irreducibility therefore gives either $P_\pi Y=0$ or $P_\pi Y=H_\pi$.

Finally, the Peter--Weyl finite-spectrum vectors are uniformly dense in every
continuous Banach representation of a compact group.  Concretely, a net of
central trigonometric-polynomial approximate identities gives operators
\[
 T_\alpha y=\int_G k_\alpha(g)U_gy\,dg
\]
that converge to $y$ in the uniform norm and are finite linear combinations
of the projections $P_\pi$.  Since $Y$ is invariant and closed, this applies
inside $Y$.  Hence
\[
 Y=\overline{\bigoplus_{\{\pi:\,P_\pi Y=H_\pi\}}H_\pi}^{\,\lVert\cdot\rVert_\infty},
\]
which is exactly the asserted universal-subcollection property.
The equivalence with $\dim V_\pi^K\leq1$ was supplied by Frobenius
reciprocity.
\end{proof}

\section{An explicit \texorpdfstring{$S_3$}{S3} witness}

Although Theorem~\ref{thm:characterization} immediately proves
Corollary~\ref{cor:s3}, the following construction makes the obstruction
fully concrete.  Put
\[
 V=\{(z_1,z_2,z_3)\in\C^3:z_1+z_2+z_3=0\},
\]
with $S_3$ permuting coordinates.  This is the irreducible standard
two-dimensional representation.  For example, with
\[
 a=\frac{(1,-1,0)}{\sqrt2},\qquad
 b=\frac{(1,1,-2)}{\sqrt6},
\]
the transpositions $(12)$ and $(23)$ have matrices
\[
 \begin{pmatrix}-1&0\\0&1\end{pmatrix},
 \qquad
 \begin{pmatrix}\tfrac12&\tfrac{\sqrt3}{2}\\
                 \tfrac{\sqrt3}{2}&-\tfrac12\end{pmatrix}.
\]
They have no common invariant line, proving irreducibility.

Equip $C(S_3)$ with the uniform-probability inner product
\[
 [f,k]=\frac1{6}\sum_{g\in S_3}f(g)\overline{k(g)},
\]
and write $\rho$ for the standard representation.  For nonzero $q\in V$,
define
\[
 (F_qv)(g)=\ip{v}{\rho(g)q},
 \qquad M_q=F_q(V)\subset C(S_3),
\]
using the convention that $\ip{\cdot}{\cdot}$ is linear in its first
argument.  For the left regular action $(U_hf)(g)=f(h^{-1}g)$,
\[
 U_hF_qv=F_q(\rho(h)v).
\]
Schur orthogonality gives, since $\dim V=2$,
\begin{equation}\label{eq:schur}
 [F_qv,F_rw]=\frac12\ip{v}{w}\ip{r}{q}.
\end{equation}
In particular, $F_q$ is injective, so $M_q$ is a closed minimal invariant
copy of $V$.

Let $c=a+b$.  Formula~\eqref{eq:schur} shows that $M_a\perp M_b$, while
\[
 F_cv=F_av+F_bv,
 \qquad
 [F_av,F_cv]=\frac12\lVert v\rVert^2\neq0
 \quad(v\neq0).
\]
Thus $M_a$ and $M_c$ are nonorthogonal.  They are distinct: if
$M_c=M_a$, then $F_bv=F_cv-F_av$ would put the nonzero space $M_b$ inside
$M_a$, contradicting $M_a\perp M_b$.  Lemma~\ref{lem:listed} says that a
universal family would have to contain both $M_a$ and $M_c$, contradicting
pairwise orthogonality.  This proves Corollary~\ref{cor:s3} directly.

\section{Verification, novelty, and scope}

The exact symbolic checker in \texttt{code/verify\_s3\_counterexample.py}
enumerates all six permutations, constructs the coefficient spaces
$M_a,M_b,M_c$, verifies their ranks and invariance, checks irreducibility via
the commutant of the generating transpositions, and confirms
$M_a\perp M_b$ but $M_a\not\perp M_c$ with $M_a\neq M_c$.  All checks pass.
The computation only verifies the displayed finite example; the proof above
is exact and does not depend on numerical evidence.

A bounded novelty search on 11 August 2026 checked the run's lightweight
indexes; the exact arXiv id and title; the author's name together with the
conjecture; the phrases ``pairwise orthogonal $G$-minimal'' and
``multiplicity-free''; arXiv:2308.09799; and the earlier conditional paper
arXiv:2110.12060.  The source remains arXiv v1, submitted 18 August 2023.
No later paper explicitly proving or disproving Conjecture~1.1 was found.
The mechanism itself is a direct application of standard compact-group
representation theory, so novelty confidence is moderate rather than high.

The packet settles the conjecture as stated and also characterizes exactly
which compact homogeneous actions satisfy it.  It does not address modified
statements that restrict the class of invariant subspaces or permit the
orthogonal family to depend on the subspace being represented.

\textbf{Human-review recommendation.}  Check first the short
minimal-space Lemma~\ref{lem:listed} and the three coefficient spaces in the
$S_3$ construction.  The general characterization can then be audited through
the standard isotypic projection and finite-spectrum density steps.

\begin{thebibliography}{9}
\bibitem{Hokamp}
S. A. Hokamp,
\emph{A Class of Homeomorphisms on Homogeneous Spaces of a Group Action},
arXiv:2308.09799, 2023.

\bibitem{PeterWeyl}
F. Peter and H. Weyl,
\emph{Die Vollst\"andigkeit der primitiven Darstellungen einer geschlossenen
kontinuierlichen Gruppe},
Math. Ann. 97 (1927), 737--755.

\bibitem{Folland}
G. B. Folland,
\emph{A Course in Abstract Harmonic Analysis},
CRC Press, 1995.
\end{thebibliography}

\end{document}
